
\documentclass[11pt]{article}
\usepackage[a4paper,margin=1in]{geometry}
\usepackage{amsmath,amssymb,booktabs,longtable,array}
\usepackage[hidelinks]{hyperref}
\usepackage{graphicx}
\usepackage{enumitem}

\title{Microneedle Superdisintegrant Film Simulator v5\\Comprehensive Physics, Mathematics, Kinetics, and Software Architecture}
\author{Technical Documentation}
\date{}

\begin{document}
\maketitle
\tableofcontents
\newpage

\section{Purpose of the Simulator}

This simulator predicts the behavior of coated microneedles containing drug-loaded HPMC films with superdisintegrants. It combines:

\begin{itemize}
\item Microneedle geometry
\item TPMS porous architectures
\item Polymer hydration
\item Capillary wicking
\item Drug release kinetics
\item Formulation effects
\item Drug loading calculations
\end{itemize}

The model is a hybrid semi-mechanistic framework. Some portions originate from classical transport theory, while others are empirical calibration relationships chosen to reproduce expected formulation trends.

\section{Simulation Workflow}

\begin{enumerate}
\item User specifies geometry.
\item User specifies formulation.
\item Material properties are calculated.
\item Geometric properties are calculated.
\item Drug loading is calculated.
\item Wicking profile is calculated.
\item Hydration kinetics are calculated.
\item Burst release kinetics are calculated.
\item Sustained release kinetics are calculated.
\item Total release profile is generated.
\item Metrics such as T50 and T80 are extracted.
\end{enumerate}

\section{Material Model}

\subsection{Inputs}

\begin{longtable}{ll}
Drug MW & Molecular weight (Da)\\
HPMC grade & E5, E15, E50, K4M\\
HPMC concentration & wt\%\\
Tween 80 concentration & wt/v\%\\
Glycerol concentration & wt\%\\
Superdisintegrant type & CCS, SSG, PVPP\\
Superdisintegrant concentration & wt\%\\
\end{longtable}

\subsection{Effective Water Diffusivity}

Implemented in:

\begin{verbatim}
effectiveDwater()
\end{verbatim}

The simulator assumes

\begin{equation}
D_w
=10^{-10}
\left(\frac{5}{\mu_{ref}}\right)^{0.28}
\exp(-0.03C_{HPMC})
\end{equation}

where

\begin{itemize}
\item $D_w$ = effective water diffusivity
\item $\mu_{ref}$ = viscosity reference for HPMC grade
\item $C_{HPMC}$ = HPMC concentration
\end{itemize}

Physical interpretation:

\begin{itemize}
\item Higher molecular entanglement reduces water transport.
\item Increasing polymer concentration reduces free volume.
\item K4M creates the strongest diffusion barrier.
\end{itemize}

\section{Viscosity Model}

The simulator estimates coating viscosity using

\begin{equation}
\mu =
\mu_{ref}
\left(
\frac{C_{HPMC}}{2}
\right)^{2.2}
\end{equation}

Examples:

\begin{center}
\begin{tabular}{ccc}
\toprule
Grade & Conc (\%) & Estimated Viscosity\\
\midrule
E5 & 5 & $\approx 37$ mPa s\\
E15 & 5 & $\approx 112$ mPa s\\
E50 & 5 & $\approx 375$ mPa s\\
K4M & 5 & $\approx 30000$ mPa s\\
\bottomrule
\end{tabular}
\end{center}

\section{Surface Tension Model}

Tween-80 lowers surface tension according to

\begin{equation}
\gamma =
72-
35\log_{10}
\left(
1+\frac{C_T}{0.01}
\right)
\end{equation}

bounded by

\begin{equation}
28 \le \gamma \le 72
\end{equation}

mN/m.

\section{Superdisintegrant Physics}

\subsection{CCS}

Mechanism:

\begin{itemize}
\item Swelling
\item Wicking
\item Fiber expansion
\end{itemize}

Default parameters:

\begin{equation}
k_0=0.14
\end{equation}

\begin{equation}
n_0=0.55
\end{equation}

\subsection{SSG}

Mechanism:

\begin{itemize}
\item Rapid swelling
\item Gel formation at high concentration
\end{itemize}

Default:

\begin{equation}
k_0=0.11
\end{equation}

\begin{equation}
n_0=0.65
\end{equation}

\subsection{PVPP}

Mechanism:

\begin{itemize}
\item Capillary action
\item Minimal gel formation
\end{itemize}

Default:

\begin{equation}
k_0=0.23
\end{equation}

\begin{equation}
n_0=0.40
\end{equation}

\section{Gel Penalty Model}

Above a critical concentration:

\begin{equation}
P_{gel}
=
0.75^{(C-C_{crit})}
\end{equation}

Consequences:

\begin{itemize}
\item Reduced hydration
\item Reduced burst release
\item Reduced overall release rate
\end{itemize}

\section{Microneedle Geometry}

\subsection{Aspect Ratio}

\begin{equation}
AR=\frac{H}{R}
\end{equation}

\subsection{Tip Angle}

\begin{equation}
\theta=\tan^{-1}\left(\frac{R}{H}\right)
\end{equation}

\section{Surface Area Models}

\subsection{Planar Film}

\begin{equation}
SA=\pi R^2
\end{equation}

\subsection{Conical Needle}

\begin{equation}
s=\sqrt{H^2+(R-r_t)^2}
\end{equation}

\begin{equation}
SA=\pi(R+r_t)s
\end{equation}

\subsection{Pyramidal Needle}

\begin{equation}
SA=4\left(\frac12 as\right)
\end{equation}

\subsection{Star Needle}

The simulator applies a 1.45 surface enhancement multiplier.

\subsection{TPMS Needle}

\begin{equation}
SA_{eff}
=SA_{external}\times SA_{factor}
\end{equation}

This approximates internal pore surface.

\section{Drug Loading Calculations}

Coating volume:

\begin{equation}
V_{coat}=SA\cdot L \cdot f_v
\end{equation}

Film density:

\begin{equation}
\rho=1.2\ g/cm^3
\end{equation}

Coating mass:

\begin{equation}
M_{coat}=V_{coat}\rho
\end{equation}

Drug mass:

\begin{equation}
M_{drug}=f_{drug}M_{coat}
\end{equation}

\section{TPMS Transport Physics}

Effective diffusivity ratio:

\begin{equation}
\frac{D_{eff}}{D}
=
\frac{\varepsilon}{\tau^2}
\end{equation}

where

\begin{itemize}
\item $\varepsilon$ = porosity
\item $\tau$ = tortuosity
\end{itemize}

Used for:

\begin{itemize}
\item Gyroid
\item Primitive
\item Diamond
\item Hybrid TPMS-core cone
\end{itemize}

\section{Lucas-Washburn Wicking}

The simulator solves

\begin{equation}
x^2=Ct
\end{equation}

or

\begin{equation}
x=\sqrt{Ct}
\end{equation}

where

\begin{equation}
C
=
\frac{r\gamma \cos\theta}
{2\eta\tau^2}
\end{equation}

Physical meaning:

\begin{itemize}
\item Larger pores increase penetration.
\item Lower viscosity increases penetration.
\item Higher surface tension increases penetration.
\item Larger tortuosity decreases penetration.
\end{itemize}

\section{Hydration Physics}

Hydration constant:

\begin{equation}
k_{hyd}
=
\frac{D_w}{L^2}\times60
\end{equation}

Film hydration:

\begin{equation}
H(t)=1-e^{-k_{hyd}t}
\end{equation}

Hydration reaches:

\begin{equation}
63.2\%
\end{equation}

at

\begin{equation}
t=\frac1{k_{hyd}}
\end{equation}

\section{Release Kinetics}

\subsection{Korsmeyer-Peppas}

\begin{equation}
F=k_{KP}t^n
\end{equation}

where

\begin{itemize}
\item $F$ = fractional release
\item $k_{KP}$ = release coefficient
\item $n$ = release exponent
\end{itemize}

\subsection{Mechanism Classification}

\begin{center}
\begin{tabular}{cc}
\toprule
n & Mechanism\\
\midrule
$\le0.45$ & Fickian\\
0.45--0.89 & Anomalous\\
$\ge0.89$ & Case II\\
\bottomrule
\end{tabular}
\end{center}

\section{Burst Release}

Burst release fraction:

\begin{equation}
F_b=B\left(1-e^{-k_bt}\right)
\end{equation}

where

\begin{equation}
k_b
=
\frac{k_{hyd}k_{burst}}
{k_{hyd}+k_{burst}}
\end{equation}

This is effectively a harmonic combination of hydration and burst kinetics.

\section{Combined Release Model}

Total release:

\begin{equation}
F_{total}
=
F_{burst}
+
(1-B)F_{sustained}
\end{equation}

where

\begin{equation}
0\le F_{total}\le1
\end{equation}

\section{Kinetic Profile Interpretation}

\subsection{T50}

Time for

\begin{equation}
F=0.5
\end{equation}

\subsection{T80}

Time for

\begin{equation}
F=0.8
\end{equation}

Fast-release systems typically achieve

\begin{equation}
T_{80}<10\ min
\end{equation}

while sustained systems exceed 20 min.

\section{TPMS Mathematical Definitions}

\subsection{Gyroid}

\begin{equation}
\sin x \cos y
+\sin y \cos z
+\sin z \cos x=0
\end{equation}

\subsection{Primitive}

\begin{equation}
\cos x+\cos y+\cos z=0
\end{equation}

\subsection{Diamond}

\begin{equation}
\sin x\sin y\sin z
+\sin x\cos y\cos z
+\cos x\sin y\cos z
+\cos x\cos y\sin z=0
\end{equation}

\section{Code Architecture}

\subsection{materialModel()}

Computes:

\begin{itemize}
\item viscosity
\item surface tension
\item diffusivity
\item wetting factor
\item burst factor
\item gel penalty
\end{itemize}

\subsection{geometryModel()}

Computes:

\begin{itemize}
\item area
\item volume
\item drug load
\item TPMS corrections
\end{itemize}

\subsection{simulateWicking()}

Computes Lucas-Washburn penetration.

\subsection{simulateRelease()}

Computes:

\begin{itemize}
\item hydration
\item burst release
\item sustained release
\item total release
\end{itemize}

\subsection{buildSim()}

Master orchestration function.

\section{Parameter Sensitivity}

Largest positive influence on release:

\begin{enumerate}
\item Higher porosity
\item Lower thickness
\item Higher CCS/PVPP content
\item Higher Tween concentration
\item Lower HPMC concentration
\end{enumerate}

Largest negative influence:

\begin{enumerate}
\item K4M
\item High HPMC loading
\item Excessive SSG
\item Large tortuosity
\item Thick coatings
\end{enumerate}

\section{Literature Basis}

Likely scientific foundations:

\begin{enumerate}
\item Korsmeyer RW et al., controlled release systems.
\item Peppas NA, polymeric drug delivery.
\item Lucas and Washburn capillary penetration theory.
\item Millington and Quirk effective diffusivity.
\item HPMC oral film literature.
\item Dissolving microneedle drug delivery literature.
\item TPMS transport and porous media theory.
\item CCS, SSG and PVPP superdisintegrant studies.
\end{enumerate}

\section{Limitations}

The simulator does not include:

\begin{itemize}
\item Finite-element diffusion
\item Skin diffusion resistance
\item Drug crystallization
\item Polymer erosion mechanics
\item Concentration-dependent diffusivity
\item Coupled swelling-deformation physics
\item Mechanical fracture of microneedles
\end{itemize}

Therefore outputs should be interpreted as engineering estimates rather than absolute predictions.

\end{document}
